\documentclass[11pt]{article}

\usepackage[a4paper,margin=1in]{geometry}
\usepackage{amsmath,amssymb,amsthm,mathtools}
\usepackage{booktabs}
\usepackage{enumitem}
\usepackage{microtype}
\usepackage[hidelinks]{hyperref}

\newtheorem{proposition}{Proposition}
\newtheorem{remark}{Remark}
\newtheorem{warning}{Important implementation warning}
\newtheorem{algorithmblock}{Algorithm}

\newcommand{\R}{\mathbb{R}}
\newcommand{\N}{\mathbb{N}}
\newcommand{\abs}[1]{\left\lvert #1\right\rvert}
\newcommand{\norm}[1]{\left\lVert #1\right\rVert}
\newcommand{\ip}[2]{\left\langle #1,#2\right\rangle}
\newcommand{\eps}{\varepsilon}
\newcommand{\supp}{\operatorname{supp}}
\newcommand{\diag}{\operatorname{diag}}

\title{Direct Certified Integration of Squared Polynomial-Zonotope Two-Jets\\
\large A drop-in optimization without constructing the squared integrand}
\author{}
\date{}

\begin{document}
\maketitle

\begin{abstract}
This note specifies a drop-in optimization for the computation of certified
\(L^2\), \(W^{1,2}\), and \(W^{2,2}\) norm enclosures from a
polynomial-zonotope two-jet.  The current implementation first constructs the
complete polynomial zonotope representing the squared norm integrand and only
then integrates it.  For the Hessian contribution, this entails a costly sparse
polynomial convolution.

The proposed routine fuses squaring, canonicalization, and integration.  It
uses symmetry twice: only the upper-triangular Hessian entries are retained,
with weight two on the off-diagonal entries, and only unordered pairs of input
monomials are processed.  More importantly, all selected two-jet coordinates
are contracted simultaneously through a weighted Gram matrix.  The complete
squared polynomial zonotope is never constructed.

The proposed method is algebraically exact: with the same pointwise-residual
semantics and the same final interval-enclosure rule, it returns the same
mathematical enclosure as the existing pipeline.  It is therefore distinct
from a remainder-norm relaxation such as an \(L^2\) triangle inequality; no
additional overestimation is introduced.
\end{abstract}

\tableofcontents

\section{Purpose and scope}

The target computation is the certified integral of one of the squared
two-jet quantities
\begin{align}
    \mathcal{S}_{0}(Y,J,H)
    &:= \norm{Y}_2^2, \label{eq:s0}\\
    \mathcal{S}_{1}(Y,J,H)
    &:= \norm{Y}_2^2+\norm{J}_F^2, \label{eq:s1}\\
    \mathcal{S}_{2}(Y,J,H)
    &:= \norm{Y}_2^2+\norm{J}_F^2+\norm{H}_F^2. \label{eq:s2}
\end{align}
For a vector-valued map \(f:\R^n\to\R^r\), the dense two-jet convention is
\[
    Y\in\R^r,\qquad
    J\in\R^{r\times n},\qquad
    H\in\R^{r\times n\times n},
\]
where \(H_{oab}=H_{oba}\).  Consequently,
\begin{equation}
    \norm{H}_F^2
    =
    \sum_{o=1}^{r}
    \left(
        \sum_{a=1}^{n}H_{oaa}^2
        2\sum_{1\leq a<b\leq n}H_{oab}^2
    \right).
    \label{eq:hessian-symmetric-square}
\end{equation}

The immediate target is the affine integration-cell path used by the adaptive
polynomial-zonotope norm routines.  The cell map is
\[
    X(\alpha)=c+G\alpha,\qquad
    \alpha\in[-1,1]^d,
\]
and its non-negative Jacobian density
\[
    \lambda:=\abs{\det G}
\]
is constant.  For an axis-aligned box, \(\lambda\) is the product of the
coordinate radii.  The physical integral is therefore
\begin{equation}
    \int_{X([-1,1]^d)}
    \mathcal{S}_{k}(Y,J,H)(x)\,dx
    =
    \lambda
    \int_{[-1,1]^d}
    \mathcal{S}_{k}(Y,J,H)(\alpha)\,d\alpha.
    \label{eq:physical-integral}
\end{equation}

The optimization described here applies to \texttt{output="interval"}.  If a
caller explicitly requests the complete squared polynomial zonotope itself,
then that object must, by definition, still be constructed.

\section{The current pipeline and its bottleneck}

Conceptually, the existing \(W^{2,2}\) path performs
\begin{equation}
\boxed{
\begin{aligned}
    &(Y,J,H)
    \\
    &\quad\longrightarrow
    Q(\eps)
    :=
    \norm{Y(\eps)}_2^2
    +\norm{J(\eps)}_F^2
    +\norm{H(\eps)}_F^2
    \\
    &\quad\longrightarrow
    \lambda Q(\eps)
    \\
    &\quad\longrightarrow
    \text{integrate domain noises and enclose the remaining uncertainty}.
\end{aligned}
}
\label{eq:old-pipeline}
\end{equation}
The expensive step is the explicit construction of \(Q\).

For one scalar polynomial-zonotope entry
\[
    z(\eps)=c+\sum_{\beta\in\mathcal{B}}a_\beta\eps^\beta,
    \qquad
    \eps^\beta:=\prod_{i=1}^{p}\eps_i^{\beta_i},
\]
the generic multiplication routine computes
\begin{align}
    z(\eps)^2
    &=
    c^2
    +2c\sum_{\beta\in\mathcal{B}}a_\beta\eps^\beta
    +\sum_{\beta,\delta\in\mathcal{B}}
        a_\beta a_\delta\eps^{\beta+\delta}.
    \label{eq:scalar-square}
\end{align}
A naive sparse product visits all \(m^2\) ordered pairs
\((\beta,\delta)\), where \(m=\abs{\mathcal{B}}\), even though
\((\beta,\delta)\) and \((\delta,\beta)\) give the same contribution.

In the motivating example, the tensor-valued Hessian has \(m=524\) distinct
exponent keys.  Squaring the three symmetry-distinct entries of a scalar-output
\(2\times2\) Hessian entails approximately
\[
    3m^2=3\cdot524^2=823{,}728
\]
ordered scalar term pairs.  The measured Hessian-squaring time was about
\(27.27\) seconds, whereas two-jet construction took about \(0.15\) seconds
and integration of the already constructed integrand took about \(0.87\)
seconds.

\section{Polynomial-zonotope noise semantics}
\label{sec:noise-semantics}

Let
\[
    \eps=(\alpha,\zeta,\eta)
\]
collect three conceptual classes of noise variables:
\begin{enumerate}[label=\arabic*.]
    \item \(\alpha\): domain variables, integrated over \([-1,1]^d\);
    \item \(\zeta\): retained symbolic approximation variables;
    \item \(\eta\): pointwise approximation-residual variables.
\end{enumerate}
The actual implementation identifies these roles from two pieces of metadata:
\texttt{noise\_kinds} and the cell's domain-noise index tuple.  The variables
need not occur in contiguous blocks.

The distinction between \(\zeta\) and \(\eta\) is essential.  A pointwise
residual symbol does not represent one fixed number over the entire domain.
Rather, it encodes an unknown function whose value may vary adversarially at
each point.  Consequently, a monomial containing a pointwise residual is not
integrated by the ordinary monomial-moment formula.

Let
\[
    Q(\eps)=q_0+\sum_{\gamma\neq0}q_\gamma\eps^\gamma
\]
be the scalar squared integrand after multiplication by neither the cell
density nor the reference measure.  Write
\[
    D:=\{\text{domain indices}\},
    \qquad
    P:=\{\text{pointwise-residual indices}\},
\]
and let \(R:=\{1,\ldots,p\}\setminus D\) be the retained indices.
For an exponent \(\gamma\), denote by \(\gamma_D\) and \(\gamma_R\) the
corresponding restrictions.

For terms containing no pointwise residual, define the box moment
\begin{equation}
    \mu_{\gamma_D}
    :=
    \int_{[-1,1]^d}\alpha^{\gamma_D}\,d\alpha
    =
    \begin{cases}
    \displaystyle
    \prod_{i\in D}\frac{2}{\gamma_i+1},
        &\gamma_i\text{ even for every }i\in D,\\[3mm]
    0,&\text{otherwise}.
    \end{cases}
    \label{eq:moment}
\end{equation}
The reference measure is
\[
    M:=2^d.
\]

The present \texttt{pointwise\_interval} integration semantics produce
\begin{align}
    b_\rho
    &:=
    \lambda
    \sum_{\substack{\gamma:\,\gamma_R=\rho\\
                    \gamma_i=0\ \forall i\in P}}
    q_\gamma\mu_{\gamma_D},
    \label{eq:retained-coefficients}\\
    R_{\mathrm{pw}}
    &:=
    \lambda M
    \sum_{\substack{\gamma:\\
                    \exists i\in P:\ \gamma_i>0}}
    \abs{q_\gamma}.
    \label{eq:pointwise-radius}
\end{align}
Here \(q_0\) may be viewed as the coefficient with zero exponent, so its
contribution is included in \(b_0\).  The retained polynomial is
\[
    B(\eps_R)=b_0+\sum_{\rho\neq0}b_\rho\eps_R^\rho.
\]
It is finally interval-enclosed as
\begin{equation}
    \boxed{
    \mathcal{I}_{\mathrm{current}}(Q)
    =
    \left[
        b_0-\sum_{\rho\neq0}\abs{b_\rho}-R_{\mathrm{pw}},
        b_0+\sum_{\rho\neq0}\abs{b_\rho}+R_{\mathrm{pw}}
    \right],
    }
    \label{eq:current-functional}
\end{equation}
with outward rounding in an implementation.

\begin{warning}
For a pointwise-residual term, the current rule uses the full reference
measure \(M\), not the possibly vanishing signed moment
\(\mu_{\gamma_D}\).  Thus an odd domain factor multiplying a pointwise
residual must not be discarded by parity.
\end{warning}

\section{First symmetry: weighted upper-triangular Hessian coordinates}

Instead of treating every scalar entry separately, collect precisely the
coordinates required by the Sobolev integrand into one vector.
For \(k\in\{0,1,2\}\), define a coordinate-selection map
\[
    \Phi_k(Y,J,H)\in\R^{N_k}
\]
as follows:
\begin{align*}
    \Phi_0(Y,J,H)
    &:=
    \bigl(Y_o\bigr)_{o},\\
    \Phi_1(Y,J,H)
    &:=
    \bigl((Y_o)_o,(J_{oa})_{o,a}\bigr),\\
    \Phi_2(Y,J,H)
    &:=
    \bigl((Y_o)_o,(J_{oa})_{o,a},
           (H_{oaa})_{o,a},(H_{oab})_{o,a<b}\bigr).
\end{align*}
Associate the diagonal weight vector \(w^{(k)}\in\R^{N_k}\) given by
\[
    w_i^{(k)}
    =
    \begin{cases}
        2,&\text{if coordinate }i\text{ represents }H_{oab},\ a<b,\\
        1,&\text{otherwise}.
    \end{cases}
\]
Let
\[
    W_k:=\diag(w^{(k)}).
\]
Then
\begin{equation}
    \mathcal{S}_k(Y,J,H)
    =
    \Phi_k(Y,J,H)^\top W_k\Phi_k(Y,J,H).
    \label{eq:weighted-coordinate-form}
\end{equation}

This realizes Hessian symmetry without duplicating \(H_{oab}\) and \(H_{oba}\).
The factor two is carried as a coordinate weight.  To reproduce the existing
implementation, the stored upper-triangular entry \(H_{oab}\), \(a<b\), is
authoritative; one should not average it with \(H_{oba}\).

\section{Second symmetry: a coefficient-level Gram contraction}
\label{sec:gram}

Let the selected coordinate vector itself be a vector-valued polynomial
zonotope
\begin{equation}
    v(\eps)
    :=
    \Phi_k(Y,J,H)(\eps)
    =
    c+\sum_{\beta\in\mathcal{B}}a_\beta\eps^\beta,
    \label{eq:vector-pz}
\end{equation}
where
\[
    c,a_\beta\in\R^{N_k}.
\]
The support \(\mathcal{B}\) is the union of exponent keys occurring anywhere
in the selected \(Y\), \(J\), and \(H\) coordinates.  A missing coefficient is
filled by zero.  This is exactly the tensor-coefficient representation already
used by the polynomial-zonotope class, viewed after coordinate selection.

Substituting \eqref{eq:vector-pz} into
\eqref{eq:weighted-coordinate-form} yields
\begin{align}
    Q(\eps)
    &:=
    v(\eps)^\top W_kv(\eps) \notag\\
    &=
    c^\top W_kc
    +2\sum_{\beta\in\mathcal{B}}
        c^\top W_ka_\beta\,\eps^\beta
    +\sum_{\beta,\delta\in\mathcal{B}}
        a_\beta^\top W_ka_\delta\,\eps^{\beta+\delta}.
    \label{eq:weighted-expansion}
\end{align}
Define
\begin{equation}
    g_\beta:=c^\top W_ka_\beta,
    \qquad
    G_{\beta\delta}:=a_\beta^\top W_ka_\delta.
    \label{eq:gram-definitions}
\end{equation}
The matrix \(G\) is symmetric.  Therefore
\begin{equation}
    Q(\eps)
    =
    c^\top W_kc
    +2\sum_{\beta\in\mathcal{B}}g_\beta\eps^\beta
    +\sum_{\beta\in\mathcal{B}}G_{\beta\beta}\eps^{2\beta}
    +2\sum_{\beta<\delta}G_{\beta\delta}\eps^{\beta+\delta},
    \label{eq:unordered-expansion}
\end{equation}
where \(<\) is any fixed ordering of the exponent keys.

Equation \eqref{eq:unordered-expansion} is the central optimization:
\begin{enumerate}[label=\arabic*.]
    \item only the union support \(\mathcal{B}\) is traversed;
    \item all selected scalar coordinates are contracted simultaneously;
    \item only unordered monomial pairs \(\beta\leq\delta\) are processed;
    \item the coefficient products \(G_{\beta\delta}\) can be computed in one
    dense matrix multiplication.
\end{enumerate}

If the coefficient vectors are stored as the rows of
\[
    A\in\R^{m\times N_k},
    \qquad
    A_{\beta,:}=a_\beta^\top,
\]
then
\begin{equation}
    G=(AW_k)A^\top,
    \qquad
    g=(AW_k)c.
    \label{eq:batched-gram}
\end{equation}
There is no need to introduce square roots of the weights.  In particular,
using \(AW_k\) avoids the unnecessary rounding discrepancy that would arise
from scaling off-diagonal Hessian coordinates by \(\sqrt{2}\).

\begin{proposition}[Equality with the explicitly squared integrand]
\label{prop:integrand-equality}
Let \(Q_{\mathrm{old}}\) be the polynomial obtained by the existing sum of
scalar squares, using upper-triangular Hessian entries with factor two.
Let \(Q_{\mathrm{direct}}\) be the polynomial defined implicitly by
\eqref{eq:unordered-expansion}.  Then
\[
    Q_{\mathrm{direct}}(\eps)=Q_{\mathrm{old}}(\eps)
\]
as formal polynomials.
\end{proposition}

\begin{proof}
By construction of \(\Phi_k\) and \(W_k\),
\[
    \Phi_k(Y,J,H)^\top W_k\Phi_k(Y,J,H)
    =
    \mathcal{S}_k(Y,J,H),
\]
including the factor two on every off-diagonal Hessian square.
Expanding the left-hand side coefficientwise gives
\eqref{eq:weighted-expansion}.  Since
\[
    a_\beta^\top W_ka_\delta
    =
    a_\delta^\top W_ka_\beta,
\]
the two ordered terms associated with \(\beta\neq\delta\) combine into the
factor two in \eqref{eq:unordered-expansion}.  This changes neither coefficients
nor exponent keys.  Hence both procedures define the same formal polynomial.
\end{proof}

\section{Fusing canonicalization with integration}
\label{sec:fused}

It would be incorrect simply to take an absolute value of every pair
contribution in \eqref{eq:unordered-expansion}.  Different monomial pairs and
different two-jet coordinates may produce the same exponent \(\gamma\), and
their coefficients must first be added.  Only after this canonicalization does
the current integration routine take absolute values.

The direct routine therefore maintains two coefficient maps:
\begin{enumerate}[label=\arabic*.]
    \item \texttt{pointwise\_coeff[\(\gamma\)]}: the canonical coefficient of
    every full exponent containing a pointwise residual;
    \item \texttt{retained\_coeff[\(\rho\)]}: the already integrated
    coefficient of every retained exponent \(\rho=\gamma_R\) for terms
    containing no pointwise residual.
\end{enumerate}
It also maintains a scalar \texttt{integrated\_center}.

For every generated pair \((\gamma,s)\), meaning the contribution
\(s\eps^\gamma\), the routing rule is:
\begin{equation}
\operatorname{route}(\gamma,s)
=
\begin{cases}
    \texttt{pointwise\_coeff[\(\gamma\)]}
        \mathrel{+}= \lambda s,
        &\exists i\in P:\gamma_i>0,\\[1mm]
    \text{discard},
        &\mu_{\gamma_D}=0,\\[1mm]
    \texttt{retained\_coeff[\(\gamma_R\)]}
        \mathrel{+}= \lambda\mu_{\gamma_D}s,
        &\text{otherwise}.
\end{cases}
\label{eq:routing}
\end{equation}
The constant \(c^\top W_kc\) is handled directly as
\[
    \texttt{integrated\_center}
    \mathrel{+}=
    \lambda M\,c^\top W_kc.
\]
Equivalently, it can be passed through
\(\operatorname{route}(0,c^\top W_kc)\).

After all contributions have been routed, set
\begin{align}
    R_{\mathrm{pw}}
    &:=
    M\sum_{\gamma}
        \abs{\texttt{pointwise\_coeff[\(\gamma\)]}},
    \label{eq:direct-pointwise-radius}\\
    b_0
    &:=
    \texttt{retained\_coeff[0]},\\
    R_{\mathrm{sym}}
    &:=
    \sum_{\rho\neq0}
        \abs{\texttt{retained\_coeff[\(\rho\)]}}.
    \label{eq:direct-symbolic-radius}
\end{align}
The final interval is
\begin{equation}
    \boxed{
    [b_0-R_{\mathrm{sym}}-R_{\mathrm{pw}},
     b_0+R_{\mathrm{sym}}+R_{\mathrm{pw}}].
    }
    \label{eq:direct-final}
\end{equation}

\begin{proposition}[Drop-in enclosure equality]
\label{prop:enclosure-equality}
In exact arithmetic, the direct algorithm
\eqref{eq:unordered-expansion}--\eqref{eq:direct-final} returns precisely
\(\mathcal{I}_{\mathrm{current}}(Q_{\mathrm{old}})\) from
\eqref{eq:current-functional}.
\end{proposition}

\begin{proof}
Proposition~\ref{prop:integrand-equality} gives equality of the formal squared
integrands.  For non-pointwise terms, moment integration is linear, so applying
\(\mu_{\gamma_D}\) before collecting equal retained exponents produces the same
\(b_\rho\) as collecting the full squared polynomial first and integrating
afterward.  For pointwise terms, the direct algorithm retains the full exponent
key and collects every contribution before taking its absolute value.
Consequently \eqref{eq:direct-pointwise-radius} equals
\eqref{eq:pointwise-radius}.  The final interval-enclosure step is identical to
\eqref{eq:current-functional}.
\end{proof}

\begin{warning}[Why some apparently faster variants are not drop-in exact]
The following changes do not necessarily reproduce the current enclosure:
\begin{enumerate}[label=\arabic*.]
    \item Taking \(\abs{s}\) for every generated pair before equal exponent
    keys have been merged loses cancellations and generally enlarges the
    pointwise radius.
    \item Integrating the \(W^{1,2}\) and Hessian contributions into two
    separate intervals and then adding the intervals loses cancellations
    between their coefficients.
    \item Using the domain moment of a pointwise-residual term is inconsistent
    with the present pointwise-residual semantics and can be unsound.
    \item Replacing the approximation part by an \(L^2\) radius and applying a
    triangle inequality is a valid alternative enclosure, but it is not the
    same enclosure.
\end{enumerate}
\end{warning}

\section{Reference algorithm}

\begin{algorithmblock}[Direct integrated squared two-jet enclosure]
\label{alg:direct}
The inputs are two-jet polynomial zonotopes \(Y,J,H\) with aligned noise
metadata, a Sobolev order \(k\in\{0,1,2\}\), an affine cell density
\(\lambda\geq0\), domain indices \(D\), and pointwise-residual indices \(P\).
Proceed as follows.
\begin{enumerate}[label=\arabic*.,leftmargin=1.8em]
    \item Select the coordinates \(v=\Phi_k(Y,J,H)\) and their weights \(w\).
    \item Form the union support \(\mathcal{B}\) in a deterministic order,
    and extract the center vector \(c\) and coefficient matrix \(A\).
    \item Compute
    \[
        G\gets(A\diag(w))A^\top,\qquad
        g\gets(A\diag(w))c.
    \]
    \item Initialize
    \[
        \texttt{pointwise\_coeff}\gets\{\},\qquad
        \texttt{retained\_coeff}
        \gets\{0:\lambda 2^{\abs D}c^\top\diag(w)c\}.
    \]
    \item For every \(\beta\in\mathcal{B}\), call
    \(\operatorname{route}(\beta,2g_\beta)\).
    \item For \(i=1,\ldots,\abs{\mathcal{B}}\), set
    \(\beta=\mathcal{B}_i\), call
    \(\operatorname{route}(2\beta,G_{ii})\), and for every \(j>i\) call
    \[
        \operatorname{route}(\beta+\mathcal{B}_j,2G_{ij}).
    \]
    \item Compute
    \[
        R_{\mathrm{pw}}
        \gets2^{\abs D}\sum_\gamma
        \abs{\texttt{pointwise\_coeff[\(\gamma\)]}},
        \quad
        b_0\gets\texttt{retained\_coeff[0]},
    \]
    and
    \[
        R_{\mathrm{sym}}
        \gets\sum_{\rho\neq0}
        \abs{\texttt{retained\_coeff[\(\rho\)]}}.
    \]
    \item Return the outward-rounded interval
    \[
        [b_0-R_{\mathrm{sym}}-R_{\mathrm{pw}},
         b_0+R_{\mathrm{sym}}+R_{\mathrm{pw}}].
    \]
\end{enumerate}
\end{algorithmblock}

\begin{algorithmblock}[Routing one implicit squared-polynomial contribution]
\label{alg:route}
For a contribution \(s\eps^\gamma\), define
\(\operatorname{route}(\gamma,s)\) as follows.
\begin{enumerate}[label=\arabic*.,leftmargin=1.8em]
    \item If \(s=0\), return.
    \item If \(\gamma_i>0\) for some \(i\in P\), perform
    \[
        \texttt{pointwise\_coeff[\(\gamma\)]}
        \mathrel{+}=\lambda s
    \]
    and return.
    \item Compute \(\mu=\operatorname{BoxMoment}(\gamma_D)\).  If
    \(\mu=0\), return.
    \item Set \(\rho=\gamma_R\) and perform
    \[
        \texttt{retained\_coeff[\(\rho\)]}
        \mathrel{+}=\lambda\mu s.
    \]
\end{enumerate}
\end{algorithmblock}

\section{Implementation blueprint for \texttt{intervalNets}}
\label{sec:implementation}

\subsection{Recommended public and private functions}

A suitable public helper is
\begin{verbatim}
integrate_pz_twojet_squared(
    jet: PZTwoJet,
    cell: PZIntegrationCell,
    integrand_kind: Literal["l2", "w12", "w22"],
) -> Interval
\end{verbatim}
with private helpers of the form
\begin{verbatim}
_twojet_weighted_coordinates(jet, integrand_kind)
_coefficient_matrix(selected_coordinates, union_support)
_route_integrated_quadratic_term(...)
\end{verbatim}

The adaptive-cell evaluator can then replace
\begin{verbatim}
integrand = _squared_twojet_integrand(jet, integrand_kind)
contribution = integrate_over_cell(integrand, cell, output="interval")
\end{verbatim}
by
\begin{verbatim}
contribution = integrate_pz_twojet_squared(
    jet, cell, integrand_kind=integrand_kind
)
\end{verbatim}
The existing integrand constructors should remain available for diagnostics,
rendering, explicit \texttt{output="pz"}, and regression comparison.

\subsection{Coordinate extraction}

For PyTorch coefficients, form one flattened coordinate vector:
\begin{enumerate}[label=\arabic*.]
    \item append all entries of \(Y\), with weights one;
    \item for \(W^{1,2}\) and \(W^{2,2}\), append all entries of \(J\), with
    weights one;
    \item for \(W^{2,2}\), append \(H_{oaa}\) with weights one and
    \(H_{oab}\), \(a<b\), with weights two.
\end{enumerate}
Apply the identical indexing operation to the center tensor and every
coefficient tensor.  This preserves the current upper-triangular Hessian
convention.

It is useful to build a globally sorted union of exponent keys.  Deterministic
sorting makes tests reproducible and prevents hash-table insertion order from
arbitrarily changing floating-point accumulation.

\subsection{Batched coefficient contractions}

Let \(A\) have one row per union-support exponent and one column per selected
coordinate.  Then, in PyTorch notation, the central contractions are
\begin{verbatim}
weighted_A = A * weights.unsqueeze(0)
gram = weighted_A @ A.T
center_cross = weighted_A @ center
center_square = torch.dot(center * weights, center)
\end{verbatim}
This replaces hundreds of thousands of small scalar tensor multiplications by
a small number of tensor kernels.  Only the upper triangle of \texttt{gram}
needs to be routed.  If memory becomes relevant for a much larger support, use
blocked Gram computation as described in Section~\ref{sec:blocked}.

\subsection{Reusing the existing final enclosure path}

To minimize semantic drift, the direct routine may construct an
\texttt{IntegratedPZResult}, not a squared input-domain PZ:
\begin{enumerate}[label=\arabic*.]
    \item create a scalar retained polynomial zonotope from
    \texttt{retained\_coeff};
    \item store \(R_{\mathrm{pw}}\) as its \texttt{interval\_radius};
    \item call the existing \texttt{interval\_enclosure()} method.
\end{enumerate}
This small post-integration PZ contains only retained non-domain exponents.  It
is not the 55,992-term squared integrand over all noises.  Reusing the existing
method also preserves the established outward-padding behavior.

\subsection{Noise metadata and alignment}

Before any contraction:
\begin{enumerate}[label=\arabic*.]
    \item verify that \(Y,J,H\) have compatible \texttt{num\_noise};
    \item merge or validate their \texttt{noise\_kinds} exactly as the existing
    PZ addition path would;
    \item obtain \(D\) from \texttt{cell.domain\_noise\_indices};
    \item obtain \(P\) using the same predicate and the same
    \texttt{POINTWISE\_RESIDUAL\_KINDS} constant as the current integrator;
    \item retain every index outside \(D\), including pointwise indices in the
    metadata, even though terms involving them are converted to a radius.
\end{enumerate}
No new classification logic should be duplicated in the norm module.

\subsection{Scope check for the cell density}

The direct formula above assumes that
\texttt{cell.jacobian\_density} is a non-negative scalar, which is exactly the
current affine-box adaptive path.  The implementation should explicitly check
this assumption and fall back to the old pipeline for a polynomial-zonotope
density.

A future extension to a polynomial density is possible by incorporating its
coefficients into the same routing calculation.  It is not necessary for the
present optimization and should not be mixed into the first implementation.

\subsection{Floating-point equality}

The method is exact as a formal-polynomial and enclosure transformation.
However, changing the summation order or replacing scalar products by a BLAS
matrix multiplication can change the last few floating-point bits.  Therefore:
\begin{enumerate}[label=\arabic*.]
    \item promise equality of the mathematical enclosure, not bitwise identity;
    \item use the same dtype and device as the input coefficients;
    \item apply the existing outward-padding functions to the final bounds;
    \item compare old and new results with a tight numerical tolerance in
    regression tests;
    \item if fully rigorous floating-point certification is required, use
    outward-rounded contractions or add a certified roundoff allowance.
\end{enumerate}

\section{Complexity and expected improvement}

Let
\begin{align*}
    m:=\abs{\mathcal{B}},
    \qquad&
    N_k:=\text{number of selected two-jet coordinates},\\
    &
    K:=\text{number of canonical accumulator keys}.
\end{align*}

If \(m_r\) is the number of scalar monomials in coordinate \(r\), the present
entrywise squaring performs approximately
\[
    \sum_{r=1}^{N_k}m_r^2
\]
ordered coefficient products, followed by construction and traversal of the
squared PZ.

The direct weighted-Gram method performs:
\begin{enumerate}[label=\arabic*.]
    \item a dense contraction of cost \(O(m^2N_k)\), executed in optimized
    tensor code;
    \item \(m(m+1)/2\) exponent-pair routing operations;
    \item \(O(K)\) final absolute-value and enclosure operations.
\end{enumerate}
The formal arithmetic count is still quadratic in \(m\), because the exact
square of a sparse polynomial is a convolution.  The speedup comes from:
\begin{enumerate}[label=\arabic*.]
    \item halving ordered pairs through \(\beta\leq\delta\);
    \item replacing separate scalar-entry products by one coefficient-vector
    contraction;
    \item moving small tensor multiplications out of Python loops;
    \item discarding odd non-pointwise domain moments immediately;
    \item collapsing non-pointwise terms directly to retained exponents;
    \item avoiding construction and a later traversal of the squared PZ.
\end{enumerate}

For \(m=524\), the union-support upper triangle contains
\[
    \frac{m(m+1)}{2}
    =
    \frac{524\cdot525}{2}
    =
    137{,}550
\]
monomial pairs.  In the motivating \(2\times2\) scalar-output Hessian case,
this replaces approximately \(823{,}728\) ordered scalar-entry pairs by one
batched weighted Gram contraction and \(137{,}550\) routing iterations.
The exact wall-clock speedup must be measured, but the dominant
\(27.27\)-second phase is the part directly targeted by this change.

\section{Blocked implementation for larger supports}
\label{sec:blocked}

The full Gram matrix requires \(O(m^2)\) memory.  For \(m=524\) in
\texttt{float64}, it occupies only about \(2.2\) MB and is unproblematic.
For much larger supports, compute upper-triangular blocks:
\[
    G_{IJ}=(A_IW_k)A_J^\top.
\]
Route a block immediately and release it before computing the next block.
For a diagonal block, process only its upper triangle.  For an off-diagonal
block with block index \(I<J\), every pair is an unordered pair and receives
the factor two.

Blocking retains exact algebraic equivalence and reduces Gram storage from
\(O(m^2)\) to \(O(b^2)\) for a chosen block size \(b\).

\section{Tests required before replacing the old path}
\label{sec:tests}

\subsection{Randomized equivalence tests}

Generate small scalar-, vector-, matrix-, and Hessian-valued PZs with mixtures
of
\[
    \texttt{domain},\qquad
    \texttt{approximation\_symbolic},\qquad
    \texttt{approximation\_pointwise}
\]
noise kinds.  For every integrand kind, compare
\begin{align*}
    \texttt{old}
    &:=
    \texttt{integrate\_over\_cell(
       \_squared\_twojet\_integrand(jet, kind), cell)},\\
    \texttt{new}
    &:=
    \texttt{integrate\_pz\_twojet\_squared(jet, cell, kind)}.
\end{align*}
Check equality within a tolerance commensurate with the dtype and outward
padding.

\subsection{Cancellation before absolute value}

Use the pointwise variable \(\eta\) and two coordinates
\[
    z_1=1+\eta,
    \qquad
    z_2=1-\eta.
\]
Then
\[
    z_1^2+z_2^2=2+2\eta^2.
\]
The two linear \(\eta\)-coefficients cancel before the pointwise radius is
formed.  A faulty pairwise-absolute-value implementation would lose this
cancellation.

\subsection{Pointwise residual multiplied by an odd domain monomial}

Use
\[
    z(\alpha,\eta)=\alpha+\eta.
\]
Its square contains \(2\alpha\eta\).  The ordinary domain moment of this term
is zero, but the current pointwise semantics assign it a nonzero interval
radius.  The direct implementation must match the old result.

\subsection{Hessian symmetry}

Construct a tensor Hessian for which \(H_{12}\) is nonzero and verify that its
contribution is exactly
\[
    H_{11}^2+2H_{12}^2+H_{22}^2.
\]
Also verify that the lower-triangular stored entry is not accidentally added a
second time.

\subsection{Zero coefficients and unequal supports}

Test different supports in \(Y,J,H\), coefficient tensors with structural
zeros, an empty support, a constant network, and a zero network.  The union
support and zero filling must reproduce scalar-entry extraction exactly.

\subsection{Regression on the motivating notebook}

On the two-hidden-layer affine-activation example, record:
\begin{enumerate}[label=\arabic*.]
    \item the old and new \(W^{2,2}\) integral intervals;
    \item the old and new norm intervals after the square root;
    \item two-jet time;
    \item direct integrated-square time;
    \item total time per cell;
    \item total adaptive runtime.
\end{enumerate}
The old result should remain available behind a diagnostic flag until the new
path has been validated on this example.

\section{Suggested staged implementation}

\begin{enumerate}[label=\textbf{Stage \arabic*.},leftmargin=2.2cm]
    \item Implement the direct routine for affine cells and
    \texttt{output="interval"} only.  Keep the old explicit integrand
    constructors unchanged.
    \item Add randomized equivalence and semantic edge-case tests.
    \item Switch the adaptive \(L^2\), \(W^{1,2}\), and \(W^{2,2}\) interval
    paths to the direct routine, with a fallback to the old path for unsupported
    cell densities.
    \item Benchmark the motivating notebook and add optional timing diagnostics.
    \item If exponent routing is still significant, replace full-Gram routing
    by blocked routing or group exponent pairs by their summed exponent.
\end{enumerate}

\section{Summary for implementation}

The requested drop-in optimization is
\[
\boxed{
\begin{aligned}
    &(Y,J,H)
    \\
    &\longrightarrow
    \text{weighted upper-triangular two-jet coefficient vectors}
    \\
    &\longrightarrow
    \text{one symmetric Gram contraction}
    \\
    &\longrightarrow
    \text{canonicalize directly into pointwise and integrated accumulators}
    \\
    &\longrightarrow
    \text{the existing final interval-enclosure rule}.
\end{aligned}
}
\]

The mathematical squared integrand is unchanged.  Hessian symmetry is encoded
by weights \(1\) on diagonal and \(2\) on upper off-diagonal entries.
Polynomial-square symmetry is encoded by processing only
\(\beta\leq\delta\), with factor two for \(\beta<\delta\).
All selected two-jet coordinates are processed together, so cancellations are
preserved across \(Y\), \(J\), and \(H\).

One limitation is unavoidable: to reproduce the current enclosure exactly,
coefficients of equal pointwise-residual monomials must still be accumulated
before their absolute values are taken.  Thus the algorithm avoids the squared
\emph{polynomial-zonotope object}, but it cannot discard all canonical
coefficient aggregation.  This is precisely what distinguishes the method
from a faster but potentially looser pairwise-radius bound.

\end{document}
